\BourbakiExerciseGroup

\begin{enumerate}
\def\labelenumi{\arabic{enumi})}
\tightlist
\item
  \begin{enumerate}
  \def\labelenumii{\alph{enumii})}
  \tightlist
  \item
    试给出连续映射 \(f: \mathbf{X} \to \mathbf{Y}\)、\(g: \mathbf{Y} \to \mathbf{Z}\) 的例子，使得 \(g \circ f\) 是逆紧的，\(f\) 不是满射，且 \(g\) 不是逆紧的。
  \end{enumerate}
\end{enumerate}

\begin{enumerate}
\def\labelenumi{\alph{enumi})}
\setcounter{enumi}{1}
\tightlist
\item
  试给出两个逆紧映射 \(f: \mathbf{X} \to \mathbf{Y}\) 与 \(g: \mathbf{X} \to \mathbf{Z}\)（其中 \(\mathbf{X}, \mathbf{Y}, \mathbf{Z}\) 是拟紧的）的例子，使得 \(x \to (f(x), g(x))\) 不是 \(\mathbf{X}\) 到 \(\mathbf{Y} \times \mathbf{Z}\) 中的逆紧映射。
\end{enumerate}

\begin{enumerate}
\def\labelenumi{\arabic{enumi})}
\setcounter{enumi}{1}
\tightlist
\item
  设 \(f: \mathbf{X} \to \mathbf{Y}\) 为逆紧映射，\(\mathbf{A}\) 为 \(\mathbf{X}\) 的子空间。证明：若 \(\mathbf{A} = \overline{f}^1(f(\mathbf{A})) \cap \overline{\mathbf{A}}\)，则限制 \(f|_{\mathbf{A}}\) 是 \(\mathbf{A}\) 到 \(f(\mathbf{A})\) 中的逆紧映射。这一条件是否必要？
\end{enumerate}

¶ 3) 设 \((\mathbf{X}_{\alpha}, f_{\alpha\beta})\) 为拟紧 Kolmogoroff 空间的一个逆系统，其中诸 \(f_{\alpha\beta}\) 是逆紧映射。证明：若诸 \(\mathbf{X}_{\alpha}\) 非空，则 \(\varprojlim \mathbf{X}_{\alpha}\) 非空（证明：在拟紧 Kolmogoroff 空间中，每个非空闭子集都含有由单个点组成的闭子集）。

\begin{enumerate}
\def\labelenumi{\arabic{enumi})}
\setcounter{enumi}{3}
\item
  试给出连续映射 \(f: \mathbf{X} \to \mathbf{Y}\)（其中 \(\mathbf{X}\) 与 \(\mathbf{Y}\) 是豪斯多夫空间）的一个例子，使得逆像 \(\overline{f}^1(\mathbf{K})\)（对每个紧子集 \(\mathbf{K}\)，\(\mathbf{Y}\) 中的）都是紧的，但该映射不是逆紧的（参见§ 9，习题 4）。
\item
  设 \(f: \mathbf{X} \to \mathbf{Y}\) 为逆紧映射。
\end{enumerate}

\begin{enumerate}
\def\labelenumi{\alph{enumi})}
\tightlist
\item
  证明：若 X 是正则的，则 \(f(\mathbf{X})\) 是 Y 的正则子空间。
\item
  证明：若 \(f(\mathbf{X})\) 是 Y 的 Lindelöf 子空间（§ 9，习题 14），则 X 是 Lindelöf 的（利用 § 5 第 4 目命题 10）。
\end{enumerate}

¶ 6) 设 X 是子空间 \(\mathbf{X}_i (i \in I)\) 的和所成的空间，\(f\) 是 X 到拓扑空间 Y 中的连续映射。对每个 \(i \in I\)，设 \(f_i: \mathbf{X}_i \to \mathbf{Y}\) 是 \(f\) 在 \(\mathbf{X}_i\) 上的限制。证明：\(f\) 是逆紧的，当且仅当每个 \(f_i\) 都是逆紧的，且族 \((f(\mathbf{X}_i))_{i \in I}\) 是局部有限的{[}利用定理 1 d){]}。

\begin{enumerate}
\def\labelenumi{\arabic{enumi})}
\setcounter{enumi}{6}
\item
  设 \(f: \mathbf{X} \to \mathbf{X}'\) 为逆紧映射，\(\mathbf{R}\)（相应地，\(\mathbf{R}'\)）为 \(\mathbf{X}\)（相应地，\(\mathbf{X}'\)）上与 \(f\) 相容的等价关系。设 \(\overline{f}: \mathbf{X} / \mathbf{R} \to \mathbf{X}' / \mathbf{R}'\) 为由 \(f\) 诱导的商空间之间的映射。设 (i) 在 \(f\) 下，关于 \(\mathbf{R}\) 饱和的每个集的像是关于 \(\mathbf{R}'\) 饱和的，且 (ii) \(\mathbf{R}'\) 是开的。证明：在这些条件下 \(\overline{f}\) 是逆紧的{[}利用定理 1 c){]}。
\item
  证明：在 § 5 第 2 目例 3 的假设下，典范映射 \(\mathbf{X} \to \mathbf{X} / \mathbf{R}\) 是逆紧的。
\end{enumerate}

¶ 9) a) 设 X、Y 为两个拓扑空间，\(f: \mathbf{X} \to \mathbf{Y}\) 为满射闭映射（但不必连续）。证明：若 \(\overline{f}^{1}(y)\) 对每个 \(y \in \mathbf{Y}\) 都是拟紧的，则 \(\overline{f}^{1}(\mathbf{K})\) 对 Y 的每个拟紧子集 K 都是拟紧的（或直接证明该结果，或利用 § 5 习题 7 与 § 9 习题 12）。

\begin{enumerate}
\def\labelenumi{\alph{enumi})}
\setcounter{enumi}{1}
\tightlist
\item
  再设 X 是正则的。证明：若 \(\overline{f}^{1}(y)\) 对每个 \(y \in Y\) 在 X 中是闭的，则 \(\overline{f}^{1}(K)\) 对 Y 的每个拟紧子集 K 在 X 中是闭的（方法同上）。又若 Y 是局部紧的，则 f 是连续的。
\end{enumerate}

\begin{enumerate}
\def\labelenumi{\arabic{enumi})}
\setcounter{enumi}{9}
\tightlist
\item
  设 X、Y 为两个拓扑空间。X 与 Y 之间的对应 (G, X, Y)（《集合论》第二章，§ 3，第 1 目）称为逆紧的，如果投影 \(\mathbf{pr}_2\) 在图像 G 上的限制是逆紧映射 \(\mathbf{G} \to \mathbf{Y}\)。
\end{enumerate}

\begin{enumerate}
\def\labelenumi{\alph{enumi})}
\item
  设 \(f: \mathbf{X} \to \mathbf{Y}\) 为映射。证明：\(f\) 是连续的，当且仅当对应 \(\overline{f}^1\)（\(\mathbf{Y}\) 与 \(\mathbf{X}\) 之间的）是逆紧的；且 \(f\) 是逆紧映射，当且仅当对应 \(\overline{f}^1\)（\(\mathbf{Y}\) 与 \(\mathbf{X}\) 之间的）与 \(f\)（\(\mathbf{X}\) 与 \(\mathbf{Y}\) 之间的）都是逆紧的。由此构造一个非连续映射 \(f\) 的例子，使对应 \(f\) 是逆紧的。
\item
  若 \((\mathbf{G}, \mathbf{X}, \mathbf{Y})\) 是逆紧对应，证明：只要 A 是 X 的闭子集，\(\mathbf{G}(\mathbf{A})\) 就是 Y 的闭子集。
\end{enumerate}

¶ 11) 设 (G, X, Y) 为逆紧对应（习题 10）。

\begin{enumerate}
\def\labelenumi{\alph{enumi})}
\tightlist
\item
  证明：若 \(\mathfrak{B}\) 是 \(X\) 上由闭集组成的滤子基，则
\end{enumerate}

\[
\bigcap_ {\mathbf {M} \in \mathfrak {B}} \mathrm{G} (\mathbf {M}) = \mathrm{G} \left(\bigcap_ {\mathbf {M} \in \mathfrak {B}} \mathbf {M}\right).
\]

\begin{enumerate}
\def\labelenumi{\alph{enumi})}
\setcounter{enumi}{1}
\tightlist
\item
  对每个 \(y \in \mathbf{Y}\) 及每个邻域 \(\mathbf{V}\)（\(\overline{\mathbf{G}}(y)\) 在 \(\mathbf{X}\) 中的），证明存在一个邻域 \(\mathbf{W}\)（\(y\) 在 \(\mathbf{Y}\) 中的），使得 \(\overline{\mathbf{G}}^{1}(\mathbf{W}) \subset \mathbf{V}\)（“根作为参数的函数的连续性”）（利用 § 9 习题 6）。
\end{enumerate}

¶ 12) a) 对应 (G, X, Y) 是逆紧的，当且仅当对每个拓扑空间 Z 及每个图像 E 在 Z × X 中为闭的对应 (E, Z, X)，图像 G \(\circ\) E 在 Z × Y 中是闭的。为证条件必要，注意 G \(\circ\) E 是 (E × Y) ∩ (Z × G) 在 \(\iota_{\mathbf{Z}} \times \mathrm{pr}_{2} : \mathbf{Z} \times \mathbf{G} \to \mathbf{Z} \times \mathbf{Y}\) 下的像。为证充分性，设 \(\delta_{\mathbf{Y}}\) 为对角映射 \(\mathbf{Y} \to \mathbf{Y} \times \mathbf{Y}\)；证明：对每个子集 \(\mathbf{F}\)（\((\mathbf{Z} \times \mathbf{Y}) \times \mathbf{X}\) 的），\(\mathbf{G} \circ \mathbf{F}\) 在映射

\[
\iota_ {\mathbf {Z}} \times \delta_ {\mathbf {Y}}: \mathbf {Z} \times \mathbf {Y} \rightarrow \mathbf {Z} \times \mathbf {Y} \times \mathbf {Y}
\]

下的逆像，是在映射 \((y, z) \to (z, y)\) 之下、集

\[
\operatorname{pr} _ {2 3} (\overline {{{\mathrm{F}}}} \cap (\mathrm{G} \times \mathrm{Z})),
\]

的像，其中 \(\mathbf{F}\) 是 \(\mathbf{F}\) 在映射 \((z, y, x) \to (x, y, z)\) 下的像。b) 由 a) 推出：两个逆紧对应的复合是逆紧对应。

\begin{enumerate}
\def\labelenumi{\alph{enumi})}
\setcounter{enumi}{2}
\tightlist
\item
  设 \((\mathbf{G},\mathbf{X},\mathbf{Y})\) 为逆紧对应。证明：若 \(\mathbf{X}\) 是豪斯多夫的，则 \(\mathbf{G}\) 在 \(\mathbf{X} \times \mathbf{Y}\) 中是闭的。
\end{enumerate}

¶ 13) 设 X 为局部紧空间，Y 为任一拓扑空间。a) 证明下列四个命题等价：

\(\alpha)\) (G, X, Y) 是 X 与 Y 之间的逆紧对应。

\(\beta)\) 对每个豪斯多夫空间 \(\mathbf{X}'\)（把 \(\mathbf{X}\) 作为子空间包含），\(\mathbf{G}\) 在 \(\mathbf{X}' \times \mathbf{Y}\) 中是闭的。

\(\gamma)\) G 在 \(\mathbf{X} \times \mathbf{Y}\) 中是闭的，且对每个 \(y \in \mathbf{Y}\)，存在邻域 \(\mathbf{V}\)（\(y\) 在 \(\mathbf{Y}\) 中的）与紧子集 \(\mathbf{K}\)（\(\mathbf{X}\) 的），使得 \((\mathbf{X} - \mathbf{K}) \times \mathbf{V}\) 不与 \(\mathbf{G}\) 相交。

δ) 存在紧空间 \(X'\)（把 X 作为子空间包含），使得 G 在 \(X' \times Y\) 中是闭的。

{[}为证 \(\alpha) \Longrightarrow \beta)\)，利用习题 12 \(a)\) 与 \(c)\)。为证 \(\delta) \Longrightarrow \alpha)\)，利用第 1 目命题 2 与第 2 目定理 1 的推论 5。{]}

\begin{enumerate}
\def\labelenumi{\alph{enumi})}
\setcounter{enumi}{1}
\tightlist
\item
  证明：若 \((\mathbf{G},\mathbf{X},\mathbf{Y})\) 是逆紧的，则：
\end{enumerate}

\(\zeta)\) 对每个紧子集 \(\mathbf{L}\)（\(\mathbf{Y}\) 的），\(\overline{\mathbf{G}}(\mathbf{L})\) 是紧集。

再证明：若 Y 是局部紧的，且 G 在 \(X \times Y\) 中是闭的并满足 \(\zeta\)，则 (G, X, Y) 是逆紧的。{[}利用 \(\alpha\) 与 \(\gamma\) 的等价性。{]}当 Y 是豪斯多夫但非局部紧的空间时，后一断言不再成立（习题 4）。

\begin{enumerate}
\def\labelenumi{\arabic{enumi})}
\setcounter{enumi}{13}
\item
  设 Y 为局部紧空间，X 为拓扑空间。证明：映射 \(f: X \rightarrow Y\) 是连续的，当且仅当对每个把 Y 作为子空间包含的紧空间 \(Y'\)，f 的图像在 \(X \times Y'\) 中是闭的。{[}利用习题 13 与 10 a){]}。
\item
  设 X 为豪斯多夫空间，A 为 X 的闭子集，f 为 X --- A 到局部紧空间 Y 中的逆紧映射，Y′ 为 Y 的单点紧化，\(\omega\) 为 Y′ 中的无穷远点。设 g 为 X 到 Y′ 中的映射，它在 X --- A 上与 f 相合，且把 A 的每个点映到 \(\omega\)。证明 g 是连续的。
\end{enumerate}

¶ 16) 设 X 为局部紧空间，R 为 X 上的等价关系，C 为 R 在 X × X 中的图像。证明：C 在 X × X 中是闭的，当且仅当对 X 的每个紧子集 K，K 关于 R 的饱和集在 X 中是闭的。（为证必要性，利用第 2 目定理 1 的推论 5）。

*17) 在局部紧空间 R 上，设 S 是把 Z 的所有点等同起来所得的等价关系。证明：S 是闭的，且 R/S 是豪斯多夫的但不是局部紧的。*

*18) 设 X 为 R 的局部紧子空间 {[}o, +∞{[}。设 S 为 X 上的等价关系，其等价类是诸集 \{x, 1/x\}（对 o \textless{} x ≤ 1）以及集 \{o\}。证明：S 是开的，S 的图像在 X × X 中是闭的，且 X/S 是紧的，但 S 不是闭的。*

¶ 19) 设 X 为局部紧 σ-紧空间，R 为 X 上其图像在 X × X 中为闭的等价关系。证明：X/R 是豪斯多夫的。{[}设 (Uₙ) 是 X 的由相对紧的开集组成的覆盖，使得 Uₙ ⊂ Uₙ₊₁。设 A、B 为两个关于 R 饱和的不相交闭集。递归地定义两个递增的、关于 R 饱和的闭集序列 (Vₙ)、(Wₙ)，使得 Vₙ ∩ Wₙ = ∅，且 Vₙ ∩ Uₙ（相应地，Wₙ ∩ Uₙ）是 A ∩ Uₙ（相应地，B ∩ Uₙ）在紧子空间 Uₙ 中的邻域。利用第 3 目命题 8 与习题 16，以及 § 9 第 2 目命题 2{]}。(*)

¶ 20) a) 设 X 为非空紧空间，R 为 X 上的豪斯多夫等价关系。证明：X/R 在空间 \(\mathfrak{P}_{0}(X)\) 中（赋有拓扑 \(\mathcal{G}_{\Theta}\)；§ 8，习题 12）是闭的，当且仅当关系 R 是开的。{[}利用 § 9 习题 13 与 § 3 习题 16 证明：若 X/R 在空间 \(\mathfrak{P}_{0}(X)\) 中（赋有拓扑 \(\mathcal{G}_{\Theta}\)）是闭的，则 X/R 上的商拓扑与由 \(\mathcal{G}_{\Theta}\) 诱导的拓扑相合，然后应用 § 5 习题 7。反之，若 R 是开的，利用 § 8 习题 29 证明 X/R 关于 \(\mathcal{G}_{\Theta}\) 是闭的{]}。

*b) 设 X 为实平面 \(\mathbb{R}^2\) 的局部紧子空间，它由诸点 \((\mathtt{I}/n, y)\)（其中 \(n\) 是整数 \(>0\) 且 \(-\mathtt{I} \leqslant y \leqslant \mathtt{I}\)）以及诸点 \((\mathtt{o}, y)\)（其中 \(-\mathtt{I} \leqslant y < 0\)，或 \(0 < y \leqslant \mathtt{I}\)，或 \(y = 2\)）组成。设 \(p\) 是 \(\mathsf{pr}_1\) 在 X 上的限制，并设 R 为与 \(p\) 相伴随的等价关系。证明：R 既非开亦非闭，但 X/R 是紧的，且在空间 \(\mathfrak{P}_0(X)\) 中（赋有拓扑 \(\mathcal{C}_{\Theta}\)）是闭的。* ¶ 21) (*) a) 设 X 为局部拟紧的可达空间（§ 9，习题 29）。设 \(\mathbf{X}_0\) 为与 X 有同一底集的离散空间，并设 \(\tilde{\mathbf{X}}_0\) 为 \(\mathbf{X}_0\) 的（紧）超滤子空间（把离散空间 \(\mathbf{X}_0\) 等同于 \(\tilde{\mathbf{X}}_0\) 的一个开子空间）（§ 9，习题 26）。在 \(\tilde{\mathbf{X}}_0\) 上，设 R {[}或 R(X){]} 为超滤子 U、B 之间的如下等价关系：“U 与 B 在 X 中的极限点所成之集相同”。证明：R 是豪斯多夫的{[}利用 § 9 习题 25 b)、§ 7 习题 7 b) 以及公理 (C') 中给出的 \(\tilde{\mathbf{X}}_0\) 的拓扑的描述{]}。紧空间 \(\mathbf{X}^c = \tilde{\mathbf{X}}_0 / \mathbf{R}\) 与 X 的本原子集所成之集（§ 7，习题 7）成一一对应。若 X 是豪斯多夫的（从而是局部紧的），则 \(\mathbf{X}^c\) 当 X 紧时可等同于 X，而当 X 非紧时可等同于 X 的单点紧化。

\begin{enumerate}
\def\labelenumi{\alph{enumi})}
\setcounter{enumi}{1}
\item
  设 Y 为 X 的闭子空间。超滤子空间 \(\tilde{Y}_{0}\) 可等同于 \(Y_{0}\) 在 \(\tilde{X}_{0}\) 中的闭包，而 \(Y^{c}\) 可等同于 \(\tilde{Y}_{0}\) 在 \(X^{c}\) 中的典范像。
\item
  证明：限制到 \(X_{0}\) 上的典范映射 \(\tilde{X}_{0} \rightarrow X^{c}\) 是 \(X_{0}\) 到子空间 \(X_{*}^{c}\)（\(X^{c}\) 的）上的双射；典范映射 \(f: X_{*}^{c} \rightarrow X\) 是连续双射；且 \(X_{*}^{c}\) 的任何紧子集在该映射下的像是 X 的闭紧子集。{[}为证 f 连续，利用 b)；还需注意：在 \(\tilde{X}_{0}\) 中，\(X_{*}^{c}\) 的一个点的逆像由 X 上具有同一个给定极限点的所有超滤子组成，且若 U 是 X 上的超滤子，则它作为 \(\tilde{X}_{0}\) 上的超滤子基在 \(\tilde{X}_{0}\) 中的极限点，就是把 U 视为 \(\tilde{X}_{0}\) 的点时所得的点{]}。
\end{enumerate}

*d) 设 \(\mathbf{Z}\) 为这样的集：它是可数个空间 \(\mathbf{Y}_n\)（每个都等同于 §8 习题 7 所定义的空间 \(\mathbf{Y}\)）与由单个点组成的集 \(\{\omega\}\) 的和。在 \(\mathbf{Z}\) 上如下定义一个拓扑：对点 \(z \in \mathbf{Y}_n\)，取它在空间 \(\mathbf{Y}_n\) 中的邻域滤子作为基本邻域系；对 \(\omega\)，取子集 \(\mathbf{V}_n\)（\(\mathbf{Y}\) 的）所成之集作为基本邻域系，其中 \(\mathbf{V}_n\) 是 \(\{\omega\}\) 与诸 \(\mathbf{Y}_m\)（\(m \geqslant n\) 者）的并。如此定义的空间 \(\mathbf{Z}\) 是可达的、拟紧的且局部拟紧的；但子空间 \(Z_*^c\)（\(Z^c\) 的）不是局部紧的。*

\begin{enumerate}
\def\labelenumi{\alph{enumi})}
\setcounter{enumi}{4}
\tightlist
\item
  设 Y 为另一个局部拟紧的可达空间，u 为 X 到 Y 中的一个连续映射，使得 Y 的每个拟紧子集在 u 下的逆像是拟紧的。把 u 视为 \(X_{*}^{c}\) 到 \(Y_{*}^{c}\) 中的映射，证明：u 是连续的，且可（唯一地）延拓为连续映射 \(\mathbf{X}^c\to \mathbf{Y}^c\)。{[}先把 \(u\) 延拓为 \(\tilde{\mathbf{X}}_0\) 到 \(\tilde{\mathbf{Y}}_0\) 中的一个连续映射，并证明该延拓与等价关系 \(\mathbf{R}(\mathbf{X}),\mathbf{R}(\mathbf{Y})\) 相容{]}。
\end{enumerate}

\(^{*}f)\) 设 X 为 § 9 习题 29 c) 所定义的拟紧、非局部拟紧的可达空间。证明：在相应的空间 \(\tilde{X}_{0}\) 中，等价关系 R(X) 不是豪斯多夫的。*

